% LedModule — LED制御モジュール
\subsection{LedModule}

\begin{apibox}{LedModule --- LED制御モジュール}
\textbf{定義ファイル:} \texttt{lib/LedModule/LedModule.h / .cpp}\\
\textbf{分類:} 出力モジュール（\texttt{updateOutput()}のみオーバーライド）\\
\textbf{対象デバイス:} 汎用LED

GPIO出力によりLEDのON/OFFを制御するシンプルなモジュール。
本アーキテクチャの最小実装例として位置づけられる。
\end{apibox}

\subsubsection{機能概要}

\begin{itemize}
    \item GPIO出力によるLED ON/OFF制御
    \item \texttt{state}フラグによるシンプルな状態管理
    \item \texttt{brightness}メンバ（将来のPWM調光用に予約）
\end{itemize}

% --- Config ---
\subsubsection{Config構造体}

\begin{funcblock}{LedConfig}
LED接続ピンを定義する。

\begin{longtable}{l l l p{5.5cm}}
\toprule
\textbf{メンバ} & \textbf{型} & \textbf{制約・既定値} & \textbf{説明} \\
\midrule
\endhead
\texttt{ledPin} & \texttt{uint8\_t} & --- & LED出力GPIOピン番号 \\
\bottomrule
\end{longtable}
\end{funcblock}

% --- Data ---
\subsubsection{Data構造体}

\begin{funcblock}{LedData}
LEDの出力状態を保持する。

\begin{longtable}{l l l p{5.5cm}}
\toprule
\textbf{メンバ} & \textbf{型} & \textbf{初期値} & \textbf{説明} \\
\midrule
\endhead
\texttt{state}      & \texttt{bool}     & \texttt{false} & \texttt{true}: LED点灯、\texttt{false}: LED消灯 \\
\texttt{brightness} & \texttt{uint8\_t} & \texttt{0}     & 輝度値（0--255）。現在の実装ではGPIOのHIGH/LOWのみで未使用 \\
\bottomrule
\end{longtable}
\end{funcblock}

% --- メソッド ---
\subsubsection{メソッド}

\begin{funcblock}{LedModule(const LedConfig\& config)}
\begin{tabularx}{\textwidth}{l X}
\toprule
\textbf{項目} & \textbf{内容} \\
\midrule
引数   & \texttt{config} --- \texttt{LedConfig}へのconst参照 \\
動作   & Configを\texttt{\_config}にコピー保持する \\
\bottomrule
\end{tabularx}
\end{funcblock}

\begin{funcblock}{bool init()}
\begin{tabularx}{\textwidth}{l X}
\toprule
\textbf{項目} & \textbf{内容} \\
\midrule
戻り値   & \texttt{bool} --- 常に\texttt{true} \\
動作     & \texttt{ledPin}を\texttt{OUTPUT}モードに設定し、LOWに初期化する \\
\bottomrule
\end{tabularx}
\end{funcblock}

\begin{funcblock}{void updateOutput(SystemData\& data)}
\begin{tabularx}{\textwidth}{l X}
\toprule
\textbf{項目} & \textbf{内容} \\
\midrule
読み取り元     & \texttt{data.led} \\
タイミング制御 & 毎ループ実行 \\
動作           & \texttt{data.led.state}に応じて\texttt{digitalWrite()}でHIGH/LOWを出力する \\
\bottomrule
\end{tabularx}
\end{funcblock}

% --- 使用例 ---
\subsubsection{使用例}

\begin{lstlisting}[style=cppstyle, caption={LedModuleの使用例}]
// ProjectConfig.h
const LedConfig LED_CONFIG = {
    .ledPin = 2,  // GPIO2: オンボードLED
};

// ロジックフェーズでの使用
void applyPattern(SystemData& data) {
    // ボタン押下でLEDをトグル
    if (data.button.justPressed) {
        data.led.state = !data.led.state;
    }
}
\end{lstlisting}
